% !TEX root = ../main.tex
\section{Ejercicio propuesto 2: pruebas diferenciales y análisis experimental}

\subsection{Corpus de prueba y protocolo de medición}
Los tres compresores se ejecutaron sobre exactamente los mismos $9$ archivos generados de forma determinista por \path{python/case_generator.py} bajo licencia CC0 (Tabla \ref{tab:corpus}). Las mediciones se realizaron en Python 3.12.8 (CPython, Windows 11 AMD64, 12 núcleos lógicos) siguiendo la Sección 10 de la guía: $1$ corrida de calentamiento y $10$ repeticiones medidas con \texttt{perf\_counter\_ns()} reportando mediana, mínimo y máximo sin lectura de disco ni impresión dentro del intervalo cronometrado. Aparte se midieron las corridas de extremo a extremo (E2E) con lectura y escritura en archivos temporales sobre los $5$ archivos grandes, y la memoria pico mediante \texttt{tracemalloc} en pasadas independientes. Todas las $42$ ejecuciones confirmaron reconstrucción idéntica y coincidencia SHA-256.

\begin{table}[H]\centering\footnotesize
\caption{Conjunto mínimo de datos utilizado, tamaño original y procedencia}\label{tab:corpus}
\begin{tabularx}{\textwidth}{@{}llrl>{\raggedright\arraybackslash}X@{}}\toprule
Archivo & Categoría & $S_o$ (B) & Licencia & Descripción y procedencia \\ \midrule
\texttt{01\_vacio.bin} & vacio & 0 & CC0 & Archivo vacío de 0 bytes \\
\texttt{02\_un\_byte.bin} & un\_byte & 1 & CC0 & Archivo de 1 byte (0x5A) \\
\texttt{03\_pocos\_bytes.bin} & pocos\_bytes & 16 & CC0 & Entrada corta de 16 bytes (sobrecarga de metadatos) \\
\texttt{04\_unico\_simbolo.bin} & unico\_simbolo & 32\,768 & CC0 & Único símbolo 0x41 repetido 32 768 veces \\
\texttt{05\_texto\_espanol.txt} & texto\_espanol & 36\,000 & CC0 & Texto natural académico en español UTF-8 (36 000 bytes) \\
\texttt{06\_patrones\_repetitivos.bin} & patrones\_repetitivos & 32\,768 & CC0 & Registros estructurados con subcadenas repetitivas (32 768 bytes) \\
\texttt{07\_uniforme\_256.bin} & uniforme & 32\,768 & CC0 & Distribución uniforme de 256 valores (128 copias c/u\, semilla 20260901) \\
\texttt{08\_aleatorio.bin} & aleatorio & 32\,768 & CC0 & Bytes pseudoaleatorios reproducibles (semilla 20260902\, 32 768 bytes) \\
\texttt{09\_imagen\_comprimida.png} & ya\_comprimido & 36\,830 & CC0 & Imagen PNG RGB 128x96 comprimida internamente con zlib nivel 9 \\
\bottomrule\end{tabularx}\end{table}


\subsection{Tamaños comprimidos, razón de compresión y metadatos}
La Tabla \ref{tab:tamanos} presenta el tamaño total $S_c$, la razón $R$ y el ahorro $A$ para los $9$ archivos, y la Tabla \ref{tab:metadatos} desglosa encabezado, metadatos y datos codificados. Las Figuras \ref{fig:ahorro} y \ref{fig:metadatos} muestran el porcentaje de ahorro y el peso relativo de los metadatos frente a los datos codificados.

\begin{table}[H]\centering\scriptsize
\caption{Tamaños totales $S_c$ (B), razón $R$ y ahorro $A$ por archivo y algoritmo (reconstrucción exacta 100\,\%)}\label{tab:tamanos}
\begin{tabular}{@{}lrrrrrrrrrr@{}}\toprule
Archivo & $S_o$ & \multicolumn{3}{c}{Huffman canónico} & \multicolumn{3}{c}{RLE binario} & \multicolumn{3}{c}{LZ77 (4 KiB)} \\
\cmidrule(lr){3-5}\cmidrule(lr){6-8}\cmidrule(l){9-11}
 & (B) & $S_c$ & $R$ & $A$ (\%) & $S_c$ & $R$ & $A$ (\%) & $S_c$ & $R$ & $A$ (\%) \\ \midrule
\texttt{01\_vacio.bin} & 0 & 17 & 0.00 & 0.0 & 18 & 0.00 & 0.0 & 22 & 0.00 & 0.0 \\
\texttt{02\_un\_byte.bin} & 1 & 20 & 0.05 & -1900.0 & 21 & 0.05 & -2000.0 & 24 & 0.04 & -2300.0 \\
\texttt{03\_pocos\_bytes.bin} & 16 & 48 & 0.33 & -200.0 & 63 & 0.25 & -293.8 & 40 & 0.40 & -150.0 \\
\texttt{04\_unico\_simbolo.bin} & 32768 & 4115 & 7.96 & 87.4 & 21 & 1560.38 & 99.9 & 3893 & 8.42 & 88.1 \\
\texttt{05\_texto\_espanol.txt} & 36000 & 20166 & 1.79 & 44.0 & 107373 & 0.34 & -198.3 & 5131 & 7.02 & 85.8 \\
\texttt{06\_patrones\_repetitivos.bin} & 32768 & 18121 & 1.81 & 44.7 & 98136 & 0.33 & -199.5 & 3976 & 8.24 & 87.9 \\
\texttt{07\_uniforme\_256.bin} & 32768 & 33297 & 0.98 & -1.6 & 97956 & 0.33 & -198.9 & 36880 & 0.89 & -12.6 \\
\texttt{08\_aleatorio.bin} & 32768 & 33297 & 0.98 & -1.6 & 97947 & 0.33 & -198.9 & 36880 & 0.89 & -12.6 \\
\texttt{09\_imagen\_comprimida.png} & 36830 & 37359 & 0.99 & -1.4 & 110061 & 0.33 & -198.8 & 41434 & 0.89 & -12.5 \\
\bottomrule\end{tabular}\end{table}

\begin{table}[H]\centering\scriptsize
\caption{Desglose real de encabezado, metadatos y datos codificados (bytes)}\label{tab:metadatos}
\setlength{\tabcolsep}{4.5pt}
\begin{tabular}{@{}llrrrrrr@{}}\toprule
Archivo & Algoritmo & $S_o$ & Encab. & Metadatos & Datos & Relleno (bits) & Total $S_c$ \\ \midrule
\texttt{03\_pocos\_bytes.bin} & Huffman & 16 & 6 & 35 & 7 & 0 & 48 \\
\texttt{03\_pocos\_bytes.bin} & RLE & 16 & 6 & 12 & 45 & 0 & 63 \\
\texttt{03\_pocos\_bytes.bin} & LZ77 & 16 & 6 & 16 & 18 & 0 & 40 \\
\texttt{04\_unico\_simbolo.bin} & Huffman & 32768 & 6 & 13 & 4096 & 0 & 4115 \\
\texttt{04\_unico\_simbolo.bin} & RLE & 32768 & 6 & 12 & 3 & 0 & 21 \\
\texttt{04\_unico\_simbolo.bin} & LZ77 & 32768 & 6 & 16 & 3871 & 0 & 3893 \\
\texttt{05\_texto\_espanol.txt} & Huffman & 36000 & 6 & 115 & 20045 & 4 & 20166 \\
\texttt{05\_texto\_espanol.txt} & RLE & 36000 & 6 & 12 & 107355 & 0 & 107373 \\
\texttt{05\_texto\_espanol.txt} & LZ77 & 36000 & 6 & 16 & 5109 & 0 & 5131 \\
\texttt{07\_uniforme\_256.bin} & Huffman & 32768 & 6 & 523 & 32768 & 0 & 33297 \\
\texttt{07\_uniforme\_256.bin} & RLE & 32768 & 6 & 12 & 97938 & 0 & 97956 \\
\texttt{07\_uniforme\_256.bin} & LZ77 & 32768 & 6 & 16 & 36858 & 0 & 36880 \\
\texttt{09\_imagen\_comprimida.png} & Huffman & 36830 & 6 & 523 & 36830 & 0 & 37359 \\
\texttt{09\_imagen\_comprimida.png} & RLE & 36830 & 6 & 12 & 110043 & 0 & 110061 \\
\texttt{09\_imagen\_comprimida.png} & LZ77 & 36830 & 6 & 16 & 41412 & 0 & 41434 \\
\bottomrule\end{tabular}\end{table}

\noindent\textbf{Lectura cuantitativa de los resultados.} En el texto en español ($S_o=36000$~B), Huffman alcanza $H(X)=4.4170$ y $L=4.4543$ bits/símbolo ($\eta=99.16$\,\%), reduciendo el archivo a $20166$~B (43.98\,\% de ahorro), mientras LZ77 aprovecha repeticiones de palabras y frases para llegar a $5131$~B (85.75\,\% de ahorro) y RLE lo expande a $107373$~B (-198.26\,\%). En la entrada homogénea de $32\,768$~B, RLE logra una razón $R=1560.38$ ($S_c=21$~B, 99.94\,\% de ahorro). En contraste, en bytes pseudoaleatorios y en el PNG ya comprimido ningún método logra ahorro: Huffman añade $+529$~B por tabla canónica y LZ77 añade $+4604$~B sobre el PNG por banderas de literales.\par


\begin{figure}[H]\centering
\includegraphics[width=.94\textwidth]{fig1_ahorro_por_archivo.png}
\caption{Porcentaje de ahorro por archivo y algoritmo considerando el tamaño total $S_c$.}\label{fig:ahorro}
\end{figure}

\begin{figure}[H]\centering
\includegraphics[width=.94\textwidth]{fig3_desglose_metadatos_datos.png}
\caption{Desglose entre encabezado+metadatos y datos codificados en escala logarítmica.}\label{fig:metadatos}
\end{figure}

\subsection{Entropía, longitud promedio y casos de expansión}
La Tabla \ref{tab:huffman_entropia} y la Figura \ref{fig:entropia} confirman $H(X) \le L < H(X)+1$ en todo el corpus. Aunque Huffman mantiene $\eta > 99$\% en archivos de alta entropía o en entradas cortas con pocos símbolos, el encabezado y la tabla canónica superan cualquier ganancia o la entropía impide comprimir (Tabla \ref{tab:expansion}).

\begin{table}[H]\centering\scriptsize
\caption{Entropía de Shannon $H(X)$, longitud promedio $L$, eficiencia $\eta$ y efecto de metadatos en Huffman canónico}\label{tab:huffman_entropia}
\resizebox{\textwidth}{!}{%
\begin{tabular}{@{}lrrrrrrrr@{}}\toprule
Archivo & Símb. & $H(X)$ & $L$ & $\eta=H/L$ & Flujo teór. (B) & Datos (B) & Enc.+Meta (B) & Total $S_c$ \\ \midrule
\texttt{01\_vacio.bin} & 0 & 0.0000 & 0.0000 & 1.0000 & 0 & 0 & 17 & 17 \\
\texttt{02\_un\_byte.bin} & 1 & 0.0000 & 1.0000 & 0.0000 & 1 & 1 & 19 & 20 \\
\texttt{03\_pocos\_bytes.bin} & 12 & 3.5000 & 3.5000 & 1.0000 & 7 & 7 & 41 & 48 \\
\texttt{04\_unico\_simbolo.bin} & 1 & 0.0000 & 1.0000 & 0.0000 & 4096 & 4096 & 19 & 4115 \\
\texttt{05\_texto\_espanol.txt} & 52 & 4.4170 & 4.4543 & 0.9916 & 20045 & 20045 & 121 & 20166 \\
\texttt{06\_patrones\_repetitivos.bin} & 31 & 4.3862 & 4.4048 & 0.9958 & 18042 & 18042 & 79 & 18121 \\
\texttt{07\_uniforme\_256.bin} & 256 & 8.0000 & 8.0000 & 1.0000 & 32768 & 32768 & 529 & 33297 \\
\texttt{08\_aleatorio.bin} & 256 & 7.9938 & 8.0000 & 0.9992 & 32768 & 32768 & 529 & 33297 \\
\texttt{09\_imagen\_comprimida.png} & 256 & 7.9876 & 8.0000 & 0.9985 & 36830 & 36830 & 529 & 37359 \\
\bottomrule\end{tabular}}%
\end{table}


\begin{figure}[H]\centering
\includegraphics[width=.94\textwidth]{fig4_entropia_vs_longitud_huffman.png}
\caption{Entropía $H(X)$, longitud promedio $L$ de Huffman y costo total real $8S_c/S_o$.}\label{fig:entropia}
\end{figure}

\begin{table}[H]\centering\footnotesize
\caption{Casos representativos donde el algoritmo expande el tamaño ($S_c > S_o$)}\label{tab:expansion}
\begin{tabularx}{\textwidth}{@{}llrrr>{\raggedright\arraybackslash}X@{}}\toprule
Archivo & Algoritmo & $S_o$ (B) & $S_c$ (B) & Incremento & Causa estructural \\ \midrule
\texttt{03\_pocos\_bytes.bin} & Huffman & 16 & 48 & +32 B & Encabezado y metadatos superan la entrada corta \\
\texttt{03\_pocos\_bytes.bin} & RLE & 16 & 63 & +47 B & Encabezado y metadatos superan la entrada corta \\
\texttt{05\_texto\_espanol.txt} & RLE & 36000 & 107373 & +71373 B & Corridas cortas (3 bytes por par exceden 1 byte original) \\
\texttt{07\_uniforme\_256.bin} & Huffman & 32768 & 33297 & +529 B & Alta entropía (~8 bits/símbolo) más tabla canónica de 256 símbolos \\
\texttt{08\_aleatorio.bin} & LZ77 & 32768 & 36880 & +4112 B & Sin coincidencias >= 3 B: sobrecarga de 1 bit de bandera por literal \\
\texttt{09\_imagen\_comprimida.png} & Huffman & 36830 & 37359 & +529 B & Alta entropía (~8 bits/símbolo) más tabla canónica de 256 símbolos \\
\texttt{09\_imagen\_comprimida.png} & RLE & 36830 & 110061 & +73231 B & Corridas cortas (3 bytes por par exceden 1 byte original) \\
\bottomrule\end{tabularx}\end{table}


\subsection{Tiempos en núcleo frente a disco y sensibilidad de LZ77}
La Tabla \ref{tab:tiempos} reporta mediana, mínimo y máximo en núcleo junto con la mediana E2E en disco y el pico de memoria; la Figura \ref{fig:tiempos} contrasta compresión y descompresión. En descompresión, RLE y LZ77 copian bloques directamente y superan en velocidad a la lectura bit a bit de Huffman en Python puro. La Tabla \ref{tab:lz77_sens} resume cómo el tamaño de ventana $W$ y $\text{max\_match}$ modifican el ahorro y la latencia de LZ77.

\begin{table}[H]\centering\scriptsize
\caption{Tiempos en núcleo (mín./med./máx., 10 corridas), mediana E2E en disco y memoria pico (\texttt{tracemalloc})}\label{tab:tiempos}
\resizebox{\textwidth}{!}{%
\begin{tabular}{@{}llrrrrrrrrr@{}}\toprule
Archivo & Algoritmo & \multicolumn{3}{c}{Compresión núcleo (ms)} & $t_c$ E2E & \multicolumn{3}{c}{Descomp. núcleo (ms)} & $t_d$ E2E & Pico comp. \\
\cmidrule(lr){3-5}\cmidrule(lr){7-9}
 & & Mín. & Med. & Máx. & (ms) & Mín. & Med. & Máx. & (ms) & (KiB) \\ \midrule
\texttt{04\_unico\_simbolo.bin} & Huffman & 8.20 & 8.28 & 11.43 & 10.08 & 12.85 & 13.18 & 19.05 & 15.89 & 10.0 \\
\texttt{04\_unico\_simbolo.bin} & RLE & 3.43 & 3.62 & 17.05 & 4.95 & 0.01 & 0.03 & 0.07 & 1.63 & 32.3 \\
\texttt{04\_unico\_simbolo.bin} & LZ77 & 25.89 & 26.47 & 27.75 & 29.61 & 5.40 & 5.52 & 17.99 & 7.56 & 28.1 \\
\texttt{05\_texto\_espanol.txt} & Huffman & 15.11 & 17.84 & 27.28 & 17.20 & 51.39 & 58.12 & 70.92 & 58.98 & 68.0 \\
\texttt{05\_texto\_espanol.txt} & RLE & 10.64 & 11.21 & 21.12 & 14.30 & 23.29 & 24.17 & 25.86 & 31.93 & 327.5 \\
\texttt{05\_texto\_espanol.txt} & LZ77 & 37.42 & 38.12 & 53.69 & 51.10 & 6.88 & 7.20 & 8.12 & 9.83 & 1745.8 \\
\texttt{06\_patrones\_repetitivos.bin} & Huffman & 12.76 & 13.57 & 19.19 & 14.85 & 43.11 & 44.32 & 60.10 & 48.56 & 51.6 \\
\texttt{06\_patrones\_repetitivos.bin} & RLE & 9.71 & 10.27 & 21.94 & 11.40 & 21.74 & 22.28 & 35.42 & 25.45 & 296.4 \\
\texttt{06\_patrones\_repetitivos.bin} & LZ77 & 36.29 & 42.34 & 73.03 & 40.94 & 6.14 & 7.35 & 12.19 & 8.34 & 363.0 \\
\texttt{08\_aleatorio.bin} & Huffman & 16.67 & 17.36 & 29.03 & 18.60 & 73.41 & 77.53 & 87.33 & 81.52 & 192.5 \\
\texttt{08\_aleatorio.bin} & RLE & 9.38 & 9.78 & 20.72 & 11.04 & 21.27 & 22.46 & 31.92 & 24.10 & 296.0 \\
\texttt{08\_aleatorio.bin} & LZ77 & 216.70 & 236.10 & 251.81 & 285.59 & 13.46 & 13.93 & 15.66 & 20.67 & 29705.4 \\
\texttt{09\_imagen\_comprimida.png} & Huffman & 18.79 & 19.69 & 24.03 & 21.69 & 81.91 & 87.01 & 94.45 & 98.17 & 202.1 \\
\texttt{09\_imagen\_comprimida.png} & RLE & 10.67 & 11.76 & 20.84 & 12.66 & 23.82 & 33.06 & 57.61 & 27.89 & 332.7 \\
\texttt{09\_imagen\_comprimida.png} & LZ77 & 229.53 & 257.72 & 272.46 & 334.76 & 15.36 & 16.19 & 23.24 & 22.66 & 33229.3 \\
\bottomrule\end{tabular}}%
\end{table}


\begin{figure}[H]\centering
\includegraphics[width=.96\textwidth]{fig2_tiempos_compresion_descompresion.png}
\caption{Tiempos medianos en núcleo (barras) y extremo a extremo con disco ($\blacktriangle$).}\label{fig:tiempos}
\end{figure}

\begin{table}[H]\centering\scriptsize
\caption{Sensibilidad de LZ77 al tamaño de ventana $W$ y longitud máxima de coincidencia}\label{tab:lz77_sens}
\setlength{\tabcolsep}{4.5pt}
\begin{tabular}{@{}lrrrrrrr@{}}\toprule
Archivo & Ventana $W$ & $\text{max\_match}$ & $S_c$ (B) & Ahorro (\%) & Literales & Coincid. & $t_c$ (ms) \\ \midrule
\texttt{05\_texto\_espanol.txt} & 64 & 18 & 36335 & -0.93 & 28328 & 2091 & 186.45 \\
\texttt{05\_texto\_espanol.txt} & 512 & 18 & 27225 & 24.38 & 13262 & 5780 & 134.94 \\
\texttt{05\_texto\_espanol.txt} & 4096 & 6 & 13332 & 62.97 & 524 & 5986 & 73.67 \\
\texttt{05\_texto\_espanol.txt} & 4096 & 18 & 5131 & 85.75 & 521 & 2128 & 43.42 \\
\texttt{06\_patrones\_repetitivos.bin} & 64 & 18 & 32246 & 1.59 & 24952 & 1954 & 129.69 \\
\texttt{06\_patrones\_repetitivos.bin} & 512 & 18 & 4802 & 85.35 & 556 & 1955 & 42.72 \\
\texttt{06\_patrones\_repetitivos.bin} & 4096 & 6 & 11691 & 64.32 & 78 & 5450 & 50.08 \\
\texttt{06\_patrones\_repetitivos.bin} & 4096 & 18 & 3976 & 87.87 & 75 & 1821 & 45.36 \\
\bottomrule\end{tabular}\end{table}

